\section{Accuracy, cost and the value of inspection}
\label{sec:inspect}
The earlier sections kept finding that the most accurate filter is not the cheapest. We finally
measured that trade-off and the value of inspection, first exactly in a simple model and then on
held-out data.

\subsection{An exact model}
\label{sec:t6}
\textbf{Code:} \code{analytic.py}, \code{experiments/t6\_analytic.py}.

To see the trade-offs of Sections~\ref{sec:sweep}--\ref{sec:shift} without any fitting, \code{analytic.py} uses a one-dimensional model: real
emails have a score $x\sim N(-1,1)$, spam has $x\sim N(+1,1)$, and $40\%$ of emails are spam. Here the
accuracy, cost and inspection rate of any threshold rule have closed forms. Table~\ref{tab:t6}
compares four rules.

\begin{table}[!htbp]
\centering\small
\input{\resultsdir/t6_rule_comparison.tex}
\caption{Exact performance in the one-dimensional model ($\cfp=10$, $\cfn=1$). FPR: share of real
emails sent to spam; FNR: share of spam delivered.}
\label{tab:t6}
\end{table}

\begin{figure}[!htbp]
\centering
\includegraphics[width=0.92\linewidth]{t6_accuracy_vs_cost.pdf}
\caption{One-dimensional model: accuracy and expected cost of the rule ``spam iff $\post\ge t$''.
Accuracy peaks at $t=0.5$, cost is lowest at $\tstar=10/11$.}
\label{fig:t6-acc}
\end{figure}

The accuracy rule is the most accurate automatic rule ($84.6\%$) but costs $0.772$ per email, more
than delivering everything ($0.400$). Figure~\ref{fig:t6-acc} shows why: between $0.5$ and $10/11$,
every step that raises accuracy also raises the cost. Moving the cut to $\tstar=10/11$ lowers accuracy
to $73.9\%$ and the cost to $0.311$, because it almost never files a real email as spam (FPR $0.9\%$).
Allowing inspection at $\ci=0.3$ cuts the cost to $0.179$ while $41.6\%$ of emails are inspected.
Figure~\ref{fig:t6} shows how this depends on the price of inspection. The best achievable cost rises
with $\ci$, and its slope at any price equals the share of emails inspected at that price. Inspection
stops entirely at $\ci=\cstar=10/11$, whatever the class separation.

\begin{figure}[!htbp]
\centering
\includegraphics[width=0.92\linewidth]{t6_risk_reject_vs_cinsp.pdf}
\caption{One-dimensional model: lowest achievable cost (left) and share of emails inspected (right)
against the inspection price $\ci$, for three class separations $\mu$. Dotted: the best cost
without inspection. Dashed: $\cstar=10/11$, where inspection stops.}
\label{fig:t6}
\end{figure}

\subsection{On held-out data}
\label{sec:c6}
\textbf{Code:} \code{experiments/task5\_inspection.py}.

We then checked these predictions with the fitted models of Section~\ref{sec:models}. For each of the 20 splits the experiment draws
$10{,}000$ emails from the one-dimensional model and uses the standard Spambase split (standardised
features). It fits the shared Gaussian model and logistic regression, tuned on validation log-loss,
and applies three rules to the same test probabilities: the accuracy cut $0.5$, the cost cut $10/11$,
and the rule with inspection at a range of prices $\ci$.

\begin{table}[!htbp]
\centering\small
\input{\resultsdir/c6_rules.tex}
\caption{Accuracy and cost per email of the two automatic rules on test data
(mean $\pm$ standard error, 20 splits, $\cfp=10$, $\cfn=1$).}
\label{tab:c6-rules}
\end{table}

\begin{table}[!htbp]
\centering\small
\input{\resultsdir/c6_savings.tex}
\caption{Inspection at $\ci=0.3$: cost of the best automatic rule, cost with inspection, the paired
saving and the share of emails inspected.}
\label{tab:c6-savings}
\end{table}

\begin{itemize}
  \item \emph{More accurate means more expensive.} Moving the cut from $0.5$ to $10/11$ lowers the
        Gaussian model's synthetic accuracy from $84.6\%$ to $73.9\%$ and cuts its cost from $0.780$ to
        $0.309$ (Table~\ref{tab:c6-rules}). Both models show the same ordering on Spambase.
  \item \emph{The fitted models match the exact model.} At $\ci=0.3$ it predicts a cost of $0.179$
        with $41.6\%$ inspected; the fitted Gaussian model gives $0.181$ and $41.7\%$, and the curves
        agree at every price (Figure~\ref{fig:c6-syn}).
  \item \emph{Inspection pays on real data.} On Spambase, logistic regression saves $0.100$ per email
        by inspecting $24.1\%$ of them, and the Gaussian model saves $0.121$ (Table~\ref{tab:c6-savings}).
        Figure~\ref{fig:c6} shows that the saving shrinks as inspection gets dearer and vanishes at
        $\ci=10/11$, where both models stop inspecting.
  \item \emph{Not an artefact of duplicates.} $9.8\%$ of Spambase test rows also appear in their
        training split. Without them inspection still cuts the cost from $0.264$ to $0.160$ (logistic)
        and from $0.291$ to $0.173$ (Gaussian).
\end{itemize}

\begin{figure}[!htbp]
\centering
\includegraphics[width=0.92\linewidth]{c6_inspection_synthetic.pdf}
\caption{Synthetic data: test cost (left) and share inspected (right) of the fitted models against
$\ci$, with the exact values for the true posterior. The lines overlap.}
\label{fig:c6-syn}
\end{figure}

\begin{figure}[!htbp]
\centering
\includegraphics[width=0.92\linewidth]{c6_inspection_spambase.pdf}
\caption{Spambase: test cost (left) and share of emails inspected (right) against the inspection
price $\ci$ (bands: $\pm1$ standard error). The vertical line marks $\cstar=10/11$.}
\label{fig:c6}
\end{figure}
